\section{Giới thiệu chung}

\subsection{Đặt vấn đề}
Trong lĩnh vực thiết kế vi mạch (IC Design) hiện đại, mức độ tích hợp của các bóng bán dẫn (transistors) trên một chip duy nhất đã có thể đạt đến con số hàng tỉ đơn vị. Sự phức tạp khổng lồ này mang lại khả năng xử lý và tính toán tính toán vượt trội, nhưng đồng thời cũng đặt ra một thách thức vô cùng to lớn ở khâu kiểm thử (Testing) sau khi chế tạo. Việc đảm bảo một vi mạch không chứa bất kỳ lỗi vật lý nào (như đứt dây, chập mạch hay lỗi kẹt cứng logic - stuck-at faults) trước khi xuất xưởng là yêu cầu sống còn, quyết định đến chất lượng và uy tín của mọi nhà sản xuất. 

Tuy nhiên, đối với các loại mạch tuần tự (Sequential Circuits) chứa số lượng lớn các phần tử nhớ (Flip-flops hoặc Registers), việc cố gắng điều khiển và quan sát trạng thái của mạch thông qua các chân Input/Output truyền thống là một bài toán gần như bất khả thi do vướng phải nội tại liên quan đến các trạng thái chìm (Internal States). Nhận thức được "nút thắt cổ chai" này, ngành công nghiệp bán dẫn đã hình thành nên lĩnh vực Thiết kế Khả kiểm (Design for Test - DFT).

Môn học EE5217 đóng vai trò định hướng và trang bị cho kỹ sư những khái niệm thực tiễn nhất của lĩnh vực kiểm thử vi mạch DFT. Dựa trên kiến thức đó, nhóm đã lựa chọn kỹ thuật Thiết kế Quét (Scan Design) làm đề tài trọng tâm để tiến hành nghiên cứu, dịch thuật tài liệu chuẩn xác và thiết kế bài báo cáo phân tích một cách toàn diện.

\subsection{Mục tiêu của Đề tài}
Bản báo cáo này được nhóm biên soạn với những mục tiêu cốt lõi sau đây:
\begin{itemize}
    \item \textbf{Nghiên cứu nguyên lý DFT:} Khảo sát sự phức tạp trong việc kiểm thử mạch tuần tự và tìm hiểu cơ chế giải quyết mang tính nền tảng thông qua kỹ thuật gài chuỗi quét (Scan Chain).
    \item \textbf{Phân tích thuật toán Quét (Scan Design):} Dịch thuật mượt mà phần kiến thức cốt lõi (dựa trên Chương 7 của giáo trình tham khảo). Không chỉ giới hạn ở việc dịch thô, nhóm tiến hành trình bày lại dưới góc nhìn mạch lạc, \textbf{có bổ sung ví dụ mở rộng từ kiến thức tự thân} liên quan đến việc tính toán thời gian (Test Time). Cấu trúc phân tích đi từ cấu hình Quét cơ bản (Full Scan) đến các kỹ thuật quy hoạch luồng tối ưu (Multiple Scans, Partial Scans...).
    \item \textbf{Trình bày chuyên nghiệp:} Xây dựng hệ thống văn bản báo cáo có tính phân tầng, áp dụng sức mạnh sắp xếp chuẩn học thuật của LaTeX để báo cáo đạt mức độ hoàn thiện cao nhất về hình thức lẫn học thuật. Đi kèm với báo cáo bản cứng là hệ thống Slide HTML/CSS tương tác phục vụ cho công tác thuyết trình sinh động.
\end{itemize}

\subsection{Bố cục của báo cáo}
Tránh việc dàn trải nội dung, cấu trúc của báo cáo được nhóm thiết kế cô đọng thành 2 chương chính nối tiếp nhau, mang tính kế thừa về mặt kỹ thuật và tập trung thẳng vào phân tích chi tiết:
\begin{itemize}
    \item \textbf{Chương 1: Giới thiệu chung.} Cung cấp thông tin nền, bức tranh tổng quan quan trọng về lĩnh vực DFT, chỉ ra nguyên do khách quan thúc đẩy sự ra đời của Scan Design và tóm lược các mục tiêu cốt lõi của bài báo cáo trong môn học EE5217.
    \item \textbf{Chương 2: Thiết kế Khả kiểm bằng Phương pháp Quét (Design for Test by Means of Scan).} Đây là "linh hồn" của toàn bộ báo cáo, đảm nhiệm quá trình giảng giải sâu sắc và tịnh tiến. Chương 2 đi từ phân tích kỹ thuật vạch luồng Quét toàn bộ (Full Scan) kết hợp máy trạng thái Adding Machine, tối ưu hoá tiết diện và thời gian bằng Quét đa luồng (Multiple Parallel Scans), nới rộng cho thiết kế cấp thanh ghi RTL (RTL Scan Strategy). Hàng loạt điểm phân tích bổ sung được đưa vào nhằm làm sáng tỏ những thách thức trong quá trình định tuyến xung clock và tối ưu chi phí (Overhead) trên IC.
\end{itemize}

\vspace{1.5cm}
\begin{flushright}
\textit{Báo cáo được hoàn thiện với tinh thần cầu thị nghiêm túc nhất, sẵn sàng thuyết phục các tiêu chuẩn đánh giá khắt khe cả về chuyên môn phần cứng lẫn hình thức báo cáo 4.0.}
\end{flushright}

